% Ken (Tongli) Su - PhD-level AI/ML Resume - 2 pages
% Website CV source. Compile with: tectonic -X compile assets/files/Ken_Su_CV.tex
\documentclass[letterpaper,10pt]{article}

\usepackage[left=0.5in,right=0.5in,top=0.4in,bottom=0.4in]{geometry}
\usepackage{enumitem}
\usepackage{titlesec}
\usepackage[hidelinks]{hyperref}
\usepackage{xcolor}
\usepackage{tabularx}
\usepackage[T1]{fontenc}
\usepackage{charter}

\definecolor{accent}{HTML}{1F3864}

\pagestyle{empty}
\setlength{\parindent}{0pt}

\titleformat{\section}{\normalsize\bfseries\color{accent}}{}{0em}{}[\vspace{-0.62em}\rule{\textwidth}{0.8pt}]
\titlespacing*{\section}{0pt}{7pt}{4pt}

\setlist[itemize]{leftmargin=1.05em, itemsep=0.5pt, topsep=1.5pt, parsep=0pt}

\newcommand{\role}[3]{%
  \textbf{#1} \textnormal{|} #2 \hfill #3\par\vspace{-1pt}
}
\newcommand{\proj}[2]{%
  \vspace{2pt}\textit{#1} \hfill #2\par\vspace{-1pt}
}
\newcommand{\edu}[4]{%
  \textbf{#1} \hfill #2\par
  \textit{#3} \hfill #4\par\vspace{2pt}
}

\begin{document}
\small

%==================== HEADER ====================
\begin{center}
  {\LARGE\bfseries Ken (Tongli) Su}\\[3pt]
  {\small +1 404-384-4141 \textperiodcentered\ \href{mailto:ken.su@emory.edu}{ken.su@emory.edu} \textperiodcentered\
  \href{https://www.linkedin.com/in/ken-su}{linkedin.com/in/ken-su} \textperiodcentered\ Atlanta, GA}
\end{center}

%==================== SUMMARY ====================
\section{Research Summary}
{\footnotesize
Computer Science PhD student at Emory University (advisor: Liang Zhao) researching \textbf{LLM agents, alignment,
and AI interpretability and safety}. My goal is to help people accomplish complex tasks with AI while retaining meaningful control.
I am particularly interested in \textbf{multi-agent collaboration and theory of mind}, personalized assistance, and when agents
should involve humans. My work includes repository retrieval for coding agents (\textbf{LARGER}, co-first author), software world
models, and cross-modal physiological signal generation (\textbf{EMBC 2026 oral}, first author).
Before the PhD I spent 2+ years as an ML engineer shipping production LLM, retrieval, and forecasting systems across manufacturing,
finance, consumer goods, and insurance. Reviewer, NeurIPS 2026.
\par}

%==================== EDUCATION ====================
\section{Education}
\edu{Emory University}{Atlanta, GA}{Ph.D. in Computer Science --- Advisor: Liang Zhao}{Aug 2025 -- Dec 2028 (expected)}
\edu{Columbia University}{New York, NY}{M.S. in Data Science, GPA: 3.95/4.00}{Dec 2023}
\edu{Emory University}{Atlanta, GA}{B.A. in Computer Science \& B.A. in Economics (Dual Major), GPA: 3.84/4.00}{May 2022}

%==================== RESEARCH ====================
\section{Research Experience \& Publications}

\role{Research Intern}{Causal Dynamics Lab}{Jul 2026 -- Aug 2026}
Research on coding agents and software world models, including predicting the consequences of changes across software systems.
\par\vspace{4pt}

\role{Graduate Research Assistant --- LLM Agents \& Retrieval}{Emory University, Advisor: Liang Zhao}{Aug 2025 -- Present}
\proj{LARGER: Lexically Anchored Repository Graph Exploration and Retrieval --- \normalfont\textbf{co-first author}}{\normalfont arXiv:2605.16352, 2026 (under review)}
\begin{itemize}
  \item Studied repository localization errors in coding agents: lexical search can miss imports, call chains, type hierarchies,
  and code--test links, motivating structural retrieval integrated into the agent's existing search workflow.
  \item Formalized \textbf{Lexically Anchored Structural Localization}: a fixed agent policy turns its own lexical observations into structural
  graph anchors and progressively exposes compact, confidence-filtered local neighborhoods --- delivering graph evidence \textit{within} lexical
  observations rather than through new agent actions, so the method drops into existing CLI agents unchanged.
  \item Decomposed the joint objective into \textbf{graph quality} and \textbf{retrieval efficiency}, proposing a divide-and-conquer approximation that
  reduces global subgraph search to lexical anchoring plus confidence-scored local expansion, and engineered the full system: multi-language
  \textbf{AST-based graph construction}, edge weighting, and community detection for graph quality; lightweight \textbf{sidecar storage},
  \textbf{dynamic lexical propagation}, and \textbf{budget-aware filtering} for retrieval efficiency.
  \item Evaluated across \textbf{4 benchmarks} (LocBench, MuLocBench, SWE-Atlas Test Writing, SWE-Atlas Codebase QA) against \textbf{7 baselines}
  (BM25, Agentless, LocAgent, SWE-agent, OpenHands, Claude Code, Codex): \textbf{+13.9 Acc@5} tuned, \textbf{+11.8 Acc@5} with fixed hyperparameters
  (robust without per-benchmark tuning), and \textbf{+12.2 Recall@5}, surpassing strong proprietary baselines on codebase QA and test writing.
\end{itemize}
\vspace{4pt}

\role{Graduate Researcher --- Cross-Modal Generative Modeling}{Emory University, Katebi \& Clifford Lab}{2025 -- Present}
\proj{Cross-Modal Signal Translation from Fetal--Maternal ECG to Fetal Doppler Waveforms --- \normalfont\textbf{first author}}{\normalfont IEEE EMBC 2026, \textbf{Oral}}
\begin{itemize}
  \item Developed cross-modal generation from fetal and maternal ECG to Doppler waveforms, studying which aspects of fetal
  cardiovascular function are recoverable from electrical signals and where reconstruction remains limited.
  \item Designed an architecture combining \textbf{dilated convolutions} (large receptive field for long-range temporal dependencies),
  \textbf{cross-modal attention} (fetal features as queries, maternal ECG as keys/values, selectively attending to maternal segments temporally
  correlated with fetal blood flow), and \textbf{self-attention} over the fetal signal.
  \item Trained on \textbf{885 synchronized fetal/maternal ECG--Doppler envelope segments from 39 pregnancies}: \textbf{PSD MSE 49.9\,$\pm$\,15.8\,dB\textsuperscript{2}}
  (\textbf{51\% below} the two-channel baseline) and heart-rate error \textbf{4.71\,$\pm$\,0.77\,bpm} (\textbf{1.5\% better}; negligible against the
  110--160\,bpm physiological range).
  \item Compared ways to incorporate maternal ECG: \textbf{cross-modal attention reduces PSD MSE by 39\%} relative to naive
  dual-channel concatenation, supporting selective use of maternal information for waveform reconstruction.
  \item Introduced a \textbf{composite loss} balancing pointwise error, derivative error, and correlation to enforce amplitude accuracy and morphological
  fidelity simultaneously.
\end{itemize}
\vspace{4pt}

\proj{A Review of Multimodal Pretraining Strategies for Physiological Wearable Sensor Foundation Models --- \normalfont\textbf{first author}}{\normalfont Under review, 2026}
\begin{itemize}
  \item Screened arXiv, Google Scholar, Semantic Scholar, PubMed, Web of Science, and IEEE Xplore (Jan 2024 -- early 2026), synthesizing \textbf{22 papers}
  into the first review organized around \textbf{multimodal pretraining strategy} as the primary lens rather than as one topic among many.
  \item Proposed a taxonomy of \textbf{five pretraining paradigms} --- contrastive, generative/reconstruction, language-grounded alignment, missing-modality
  learning, and cross-modal knowledge distillation --- comparing objective design and architectural patterns across PPG, ECG, EEG, and ACC.
  \item Analyzed the recent model cohort (SleepFM, LSM/LSM-2, MoCA, NormWear, SensorLM, PhysioOmni, CSFM, AnyPPG) to extract shared successful principles
  (shared cross-modal representation spaces, physiology-aware positive pairing, masked cross-modal reconstruction) and open challenges in personalization,
  robustness to missing sensors, and demographic bias.
\end{itemize}

\newpage

%==================== INDUSTRY ====================
\section{Industry Experience}

\role{Machine Learning Engineer, GenAI}{DuPont Specialty Products USA, LLC}{New York, NY \textbar\ Jul 2024 -- Aug 2025}
\begin{itemize}
  \item Collaborated cross-functionally to build a scalable lead identification and prioritization pipeline on \textbf{Azure Databricks} using
  \textbf{Spark}, \textbf{Delta Lake}, and generative models, powering an internal co-pilot for sales managers that became the company's primary method
  for identifying potential customers and contributed \textbf{over \$60 million} in estimated business value.
  \item Fine-tuned a \textbf{Sentence-BERT} model with \textbf{triplet loss} on a retrieval-and-ranking task matching leads to relevant internal
  applications: \textbf{Recall@5 of 0.86} and \textbf{MRR of 0.64}, a \textbf{27\% improvement} over the existing GPT-based system, with hallucinations
  eliminated. Enabled marketing efforts that converted leads into \textbf{over \$3 million} in revenue.
  \item Engineered a data enrichment module on the \textbf{OpenAI API} with round-robin load balancing, asynchronous processing, and batch handling:
  pipeline runtime \textbf{10\,h\,$\rightarrow$\,6\,h (40\% reduction)}, token costs down \textbf{\$12,000/year}, and plug-and-play deployment saving
  \textbf{60 hours} of development time on related projects.
  \item Translated business feedback into model improvements and new features over a six-week sprint, \textbf{increasing DAU 48\%} while raising average
  user satisfaction from 4.0 to 4.6/5. Presented results to company leadership, securing additional resources for continued development.
\end{itemize}
\vspace{4pt}

\role{Machine Learning Engineer}{ICBC Financial Service, LLC}{New York, NY \textbar\ Feb 2024 -- May 2024}
\begin{itemize}
  \item Collaborated with research scientists to develop an in-house \textbf{foundation model} based on \textbf{transformers}, establishing a robust training
  loop with the Bing Search API for periodic continuous re-pretraining. \textbf{Few-shot learning} for financial report translation reached \textbf{BLEU 0.59}
  (cumulative to 4-grams), reducing external translation costs \textbf{\$30,000/year} and turnaround time \textbf{50\%}.
  \item Created a \textbf{Flask} application with Python-embedded \textbf{T-SQL} and \textbf{parameterized queries}, reducing data retrieval time \textbf{40\%}
  (saving the team 20 hours/week on compliance requests) and eliminating SQL injection risk; \textbf{85\% user adoption} among 50 team members in month one.
  \item Optimized \textbf{MySQL} scripts by rewriting queries and implementing indexing, cutting job runtime \textbf{50\%} (4\,h\,$\rightarrow$\,2\,h) and server
  load \textbf{30\%}, enabling timely investment decisions that contributed to a \textbf{3\% increase} in portfolio returns.
\end{itemize}
\vspace{4pt}

\role{Machine Learning Scientist}{Unilever United States, Inc}{New York, NY \textbar\ Sep 2023 -- Dec 2023}
\begin{itemize}
  \item Developed a \textbf{multi-modal neural network} with \textbf{PyTorch} and \textbf{Keras} combining tabular product features with image embeddings
  extracted from a fine-tuned \textbf{VGG16} model, improving \textbf{R\textsuperscript{2} by 54\%} over the baseline \textbf{XGBoost} model in sales volume
  prediction and optimizing \textbf{\$3 million} in supply chain and inventory management cost.
  \item Engineered an \textbf{ETL pipeline} for e-commerce merchant datasets ($\sim$100\,GB) in heterogeneous formats with \textbf{PySpark} and \textbf{Spark SQL},
  enhancing data quality through advanced cross-joining techniques and Levenshtein distance matching.
  \item Deployed the model as an end-to-end pipeline on \textbf{GCP} enabling scalable real-time sales forecasting with scheduled and event-trigger workflows
  and managed task dependencies, enhancing operational efficiency \textbf{20\%} per business team feedback.
\end{itemize}
\vspace{4pt}

\role{Data Scientist, Personal Insurance R\&D}{The Travelers Companies, Inc}{Hartford, CT \textbar\ Jun 2023 -- Aug 2023}
\begin{itemize}
  \item Crafted and deployed a risk assessment model using \textbf{LightGBM}, augmented with rule-based triggers and fine-tuned via \textbf{Optuna} with
  customized objective functions, achieving estimated annual savings of \textbf{\$15 million} in claim resolving cost in backtesting.
  \item Developed a \textbf{Cox Proportional Hazard} model incorporating \textbf{SMOTE} to predict occurrences of roof discoloration and damage, achieving
  \textbf{Harrell's c-index of 0.73} and enabling an inspection schedule that cut average customer life-time on-site inspection cost \textbf{12\%} per policy.
  \item Conducted extensive \textbf{exploratory data analysis} and presented findings to leadership, demonstrating model-update benefits that boosted the
  project's prioritization and lifted the referral action capture rate \textbf{33\%} (from 32.4\% to 43.1\%).
\end{itemize}

%==================== SKILLS ====================
\section{Technical Skills}
\vspace{1pt}
{\footnotesize
\begin{tabularx}{\textwidth}{@{}p{1.5in}X@{}}
  \textbf{LLM \& Generative AI} & Transformers (Hugging Face), LLM \& CLI coding agents, retrieval and ranking (RAG, BM25, graph retrieval), fine-tuning (LoRA, triplet loss), Sentence-BERT, few-shot learning, prompt design, OpenAI API, evaluation (Acc@k, Recall@k, MRR, BLEU) \\[1.5pt]
  \textbf{ML \& Deep Learning} & PyTorch, TensorFlow, Keras, Scikit-learn, LightGBM, XGBoost, Optuna; attention mechanisms, cross-modal \& multimodal architectures, generative models, self-supervised \& contrastive learning, foundation model pretraining, time-series \& biosignal modeling, survival analysis \\[1.5pt]
  \textbf{Programming Languages} & Python, SQL, R, Java, JavaScript, C++, MATLAB, Bash, HTML, STATA \\[1.5pt]
  \textbf{Big Data \& Engineering} & Spark, PySpark, Dask, Pandas, NumPy, Statsmodels, Hadoop, Airflow, Databricks, Delta Lake, MySQL, PostgreSQL, T-SQL, MongoDB, Flask, BeautifulSoup \\[1.5pt]
  \textbf{Cloud \& Visualization} & AWS (EC2, S3, Redshift), Azure (VM, DevOps), GCP (VM, Cloud Storage, BigQuery); Tableau, Power BI, Seaborn, Matplotlib \\
\end{tabularx}
}

%==================== SERVICE ====================
\section{Academic Service}
\begin{itemize}
  \item \textbf{Reviewer}, Conference on Neural Information Processing Systems (\textbf{NeurIPS}) \textbf{2026}.
  \item \textbf{Presenter (Oral)}, IEEE Engineering in Medicine \& Biology Conference (EMBC) 2026.
\end{itemize}

\end{document}
